\documentclass[12pt]{article}
\usepackage[utf8]{inputenc}
\usepackage[T1]{fontenc}
\usepackage{hyperref}
\title{Vitu Language Resources for K-12 Curriculum Development}
\author{Compiled by Claude Sonnet 5.5 (claude-sonnet-5-5) for Andr\'as Kornai}
\date{}
\begin{document}
\maketitle

\section{Introduction}
Vitu (ISO 639-3 wiv) is an Oceanic (Austronesian) language spoken on the Vitu (Witu) Islands off the northwest coast of West New Britain, Papua New Guinea. Search results gave about 7,000 speakers, against about 8,800 in the task description. It is closely related to the neighbouring Uneapa (Bali) language, and the two are sometimes treated as one, Bali-Vitu. It must not be confused with Wiru (wiu), a Southern Highlands language.

\section{Primary Educational Resources}
One search result says the language is taught in primary schools and that its orthography has been standardized. I did not see any specific Vitu primers or textbooks. Omniglot has an overview of the writing system: \url{https://www.omniglot.com/writing/vitu.htm}.

\section{Linguistic and Reference Materials}
The main grammar is \emph{Vitu Grammar Sketch} by Ren\'e van den Berg and Peter Bachet (2006, SIL, Data Papers on Papua New Guinea Languages 51). It includes two interlinearised folktales. A review is at \url{https://dev.linguistlist.org/issues/18/371/}. I did not see a direct URL for the grammar itself, so it is not in the JSON. I found no dedicated Vitu dictionary. Verbix has a short verb list of unclear source: \url{https://docs.verbix.com/Languages/Vitu}.

Other references: \url{https://en.wikipedia.org/wiki/Vitu_language} and the Glottolog record \url{https://glottolog.org/resource/languoid/id/mudu1242}.

\section{Religious and Community Resources}
My searches found none.

\section{Digital and Audio Resources}
My searches found no audio or digital Vitu resources.

\section{Curriculum Development Context}
The documentation is limited to a grammar sketch with two texts. A curriculum would need new community-based materials, building on the standardized orthography and the existing primary-school teaching of the language.
\end{document}